% from ega4-17-fr.tex
\oldpage[IV-4]{102}

  \item[b)] Les $(g_i)_x$ tels que $1\leq i\leq r$ engendrent $\mathfrak m_x$ (et par suite forment un système régulier de paramètres pour $\sh{O}_{X,x}$ (0, 17.1.6)) et les images dans $\Omega^1_{k(x)/k}$ des éléments $(d_{X/k}g_i)_x$ pour $r+1\leq i\leq n$ engendrent $\Omega^1_{k(x)/k}$.
\end{enumerate}
\end{proposition}

On a en effet (0, 20.5.12.1) la suite exacte de $k(x)$-modules
\[
\label{IV.17.15.8.1}
\tag{17.15.8.1}
\mathfrak m_x/\mathfrak m_x^2\to(\Omega^1_{X/k})_x\otimes_{\sh{O}_{X,x}}k(x)\to\Omega^1_{k(x)/k}\to0
\]
et la condition b) entraîne par suite que les $(d_{X/k}g_i)_x$ engendrent $(\Omega^1_{X/k})_x$ compte tenu du lemme de Nakayama; le fait que b) implique a) résulte donc de (17.15.3). Inversement, si a) est vérifiée, les $(d_{X/k}g_i)_x$ pour $1\leq i\leq n$ forment une base de $(\Omega^1_{X/k})_x$ en vertu de (17.15.3). Si $t_i$ ($1\leq i\leq n$) est l'image de $(g_i)_x$ dans $k(x)$, on conclut de ce qui précède que les $dt_i$ engendrent $\Omega^1_{k(x)/k}$ pour $1\leq i\leq n$, et comme par hypothèse $t_i=0$ pour $1\leq i\leq r$, les $dt_i$ pour $r+1\leq i\leq n$ engendrent déjà $\Omega^1_{k(x)/k}$. Comme $\deg.\operatorname{tr}_k k(x)=n-r$, il suit de l'égalité de Cartier (0, 21.7.1) que $\Upsilon_{k(x)/k}=0$, donc $k(x)$ est une extension séparable de $k$ (0, 20.6.3), et la suite de $k(x)$-espaces vectoriels
\[
\label{IV.17.15.8.2}
\tag{17.15.8.2}
0\to\mathfrak m_x/\mathfrak m_x^2\to(\Omega^1_{X/k})_x\otimes_{\sh{O}_{X,x}}k(x)\to\Omega^1_{k(x)/k}\to0
\]
est exacte ((0, 20.5.14) et (0, 19.6.1)); en outre les $dt_i$ pour $r+1\leq i\leq n$ forment une base de $\Omega^1_{k(x)/k}$, donc aucun des $t_i$ tels que $r+1\leq i\leq n$ ne peut être nul. Cela montre que les images des $(g_i)_x$ dans $\mathfrak m_x/\mathfrak m_x^2$ pour $1\leq i\leq r$ engendrent nécessairement $\mathfrak m_x/\mathfrak m_x^2$, donc les $(g_i)_x$ pour $1\leq i\leq r$ engendrent $\mathfrak m_x$ en vertu du lemme de Nakayama, ce qui achève de prouver que a) implique b).

\begin{corollary}[17.15.9]
\label{IV.17.15.9}
Soient $X$ un préschéma localement de type fini sur un corps $k$, $x$ un point de $X$ et posons $n=\dim_x(X)$, $r=\dim(\sh{O}_{X,x})$. Les conditions suivantes sont équivalentes:
\begin{enumerate}
  \item[a)] L'anneau $\sh{O}_{X,x}$ est régulier et $k(x)$ est une extension séparable de $k$.
  \item[b)] $X$ est lisse sur $k$ au point $x$, et l'homomorphisme canonique
  \[
  \mathfrak m_x/\mathfrak m_x^2\to(\Omega^1_{X/k})_x\otimes_{\sh{O}_{X,x}}k(x)
  \]
  est injectif.
  \item[c)] Il existe $n$ sections $g_i$ de $\sh{O}_X$ au-dessus d'un voisinage ouvert $U$ de $x$ telles que $(g_i)_x\in\mathfrak m_x$ pour $1\leq i\leq r$ et que le morphisme $U\to\operatorname{Spec}(k[T_1,\ldots,T_n])$ correspondant aux $g_i$ (cf. (17.15.8)) soit étale au point $x$.
  \item[d)] Il existe $n$ sections $g_i$ ($1\leq i\leq n$) de $\sh{O}_X$ au-dessus d'un voisinage ouvert $U$ de $x$, telles que les $(g_i)_x$ pour $1\leq i\leq r$ engendrent $\mathfrak m_x$ et que les images dans $\Omega^1_{k(x)/k}$ des $(d_{U/k}g_i)_x$ pour $r+1\leq i\leq n$ engendrent $\Omega^1_{k(x)/k}$.
\end{enumerate}
\end{corollary}

L'équivalence de c) et d) résulte de (17.15.8). De plus, on a vu dans la démonstration de (17.15.8) que la condition c) entraîne que $X$ est lisse sur $k$ au point $x$ et que la suite (17.15.8.2) est exacte, donc c) entraîne b). La condition b) entraîne que l'anneau $\sh{O}_{X,x}$ est régulier, donc $\mathfrak m_x/\mathfrak m_x^2$ est de rang $r$; en outre, $(\Omega^1_{X/k})_x$ est alors un $\sh{O}_{X,x}$-module libre de rang $n$ (17.10.2); comme la suite (17.15.7.2) est exacte par hypothèse, $\Omega^1_{k(x)/k}$ est de rang $n-r=\deg.\operatorname{tr}_k k(x)$ et l'égalité de Cartier montre que $\Upsilon_{k(x)/k}=0$, donc $k(x)$ est séparable sur $k$ (0, 20.6.3); ainsi b) entraîne a). Enfin, si a)

